\subsection{CANCELLABLE ELEMENTS}

DEFINITION 4. Given a law of composition \(\top\) on a set E, the mapping \(x \mapsto a \top x\) (resp. \(x \mapsto x \top a\) ) of E into itself is called left translation (resp. right translation) by an element \(a \in E\) .

On passing to the opposite law, left translations become right translations and conversely.

Let \(\pmb{\gamma}_{a},\pmb{\delta}_{a}\) (or \(\pmb {\gamma}(a),\pmb {\delta}(a)\) ) denote the left and right translations by \(a\in \mathbf{E}\) ; then

\[
\boldsymbol {\gamma} _ {a} (x) = a \top x, \quad \boldsymbol {\delta} _ {a} (x) = x \top a.
\]

PROPOSITION 1. If the law \(\top\) is associative, then for all \(x \in E\) and \(y \in E\)

\[
\gamma_ {x \top y} = \gamma_ {x} \circ \gamma_ {y}, \quad \delta_ {x \top y} = \delta_ {y} \circ \delta_ {x}.
\]

For all \(z \in \mathbf{E}\) :

\[
\boldsymbol {\Upsilon} _ {x \top y} (z) = (x \top y) \top z = x \top (y \top z) = \boldsymbol {\Upsilon} _ {x} (\boldsymbol {\Upsilon} _ {y} (z))
\]

\[
\boldsymbol {\delta} _ {x \top y} (z) = z \top (x \top y) = (z \top x) \top y = \boldsymbol {\delta} _ {y} (\boldsymbol {\delta} _ {x} (z))
\]

In other words, the mapping \(x \mapsto \gamma_x\) is a homomorphism from the magma E to the set \(\mathbf{E}^{\mathbb{E}}\) of mappings of E into itself with the law \((f, g) \mapsto f \circ g\) ; the mapping \(x \mapsto \delta_x\) is a homomorphism of E into the set \(\mathbf{E}^{\mathbb{E}}\) with the opposite law. If E is a monoid, these homomorphisms are unital.

DEFINITION 5. An element a of a magma E is called left (resp. right) cancellable (or regular) if left (resp. right) translation by a is injective. A left and right cancellable element is called a cancellable (or regular) element.

In other words, for \(a\) to be cancellable under the law \(\top\) , it is necessary and sufficient that each of the relations \(a \top x = a \top y\) , \(x \top a = y \top a\) imply \(x = y\) (it is said that ``a can be cancelled'' from each of these equalities). If there exists an identity element \(e\) under the law \(\top\) , it is cancellable under this law: the translations \(\gamma_e\) and \(\delta_e\) are then the identity mapping of E onto itself.

Examples. (1) Every natural number is cancellable under addition; every natural number \(\neq 0\) is cancellable under multiplication.

\begin{enumerate}
\def\labelenumi{(\arabic{enumi})}
\setcounter{enumi}{1}
\tightlist
\item
  In a lattice there can be no cancellable element under the law sup other than the identity element (least element) if it exists; similarly for inf. In particular, in the set of subsets of a set E, \(\varnothing\) is the only cancellable element under the law \(\cup\) and E the only cancellable element under the law \(\cap\) .
\end{enumerate}

PROPOSITION 2. The set of cancellable (resp. left cancellable, resp. right cancellable) elements of an associative magma is a submagma.

If \(\gamma_{x}\) and \(\gamma_{y}\) are injective so is \(\gamma_{x\top y} = \gamma_{x}\circ \gamma_{y}\) (Proposition 1). Similarly for \(\delta_{x\top y}\) .

